\documentclass[11pt,a4paper]{article}
\usepackage[utf8]{inputenc}
\usepackage[T1]{fontenc}
\usepackage{amsmath,amssymb,amsthm}
\usepackage{algorithm}
\usepackage{algpseudocode}
\usepackage{graphicx}
\usepackage{booktabs}
\usepackage{multirow}
\usepackage{hyperref}
\usepackage{xcolor}
\usepackage{enumitem}
\usepackage[margin=1in]{geometry}
\usepackage{bbm}

\newtheorem{theorem}{Theorem}
\newtheorem{proposition}{Proposition}
\newtheorem{definition}{Definition}
\newtheorem{corollary}{Corollary}

\title{HyP-DLM: Hypergraph Propagation with Dynamic Logic Modulation\\
\large DAG-Guided Multi-Hop Retrieval-Augmented Generation via Semantic Chunking and Matrix Propagation}

\author{Hieu Nguyen\\
\texttt{hieu@example.com}}

\date{}

\begin{document}
\maketitle

%% ============================================================
%% ABSTRACT
%% ============================================================
\begin{abstract}
Standard Retrieval-Augmented Generation (RAG) relies on chunk-based retrieval, while graph-based RAG methods advance this by structuring knowledge as graphs. However, existing graph-based RAG approaches are constrained by either binary relational graphs (limiting expressiveness) or expensive LLM-based knowledge extraction (limiting scalability). We propose HyP-DLM, a training-free hypergraph-based RAG framework that addresses both limitations. Our approach makes three key contributions: (1) a cost-efficient hypergraph construction pipeline using coreference resolution, semantic chunking, and local NER models (GLiNER2) instead of LLM-based extraction, where each semantic chunk naturally serves as a hyperedge connecting multiple entities; (2) a DAG-guided propagation mechanism that decomposes complex queries into directed acyclic sub-question graphs and executes degree-normalized matrix propagation with dynamic semantic modulation and confidence-weighted synonym bridging; and (3) a multi-granular scoring mechanism inspired by MaGiX that combines hyperedge-level and entity-level scores for robust passage ranking. Experiments on HotpotQA, 2WikiMultiHopQA, and MuSiQue demonstrate that HyP-DLM achieves competitive or superior performance compared to state-of-the-art graph-based RAG methods while reducing indexing cost to near-zero.
\end{abstract}

%% ============================================================
%% 1. INTRODUCTION
%% ============================================================
\section{Introduction}

Retrieval-Augmented Generation (RAG)~\cite{lewis2020rag} has become a standard paradigm for knowledge-intensive tasks, integrating knowledge retrieval with large language models (LLMs) to improve factual accuracy. Standard RAG segments documents into fixed-length chunks and retrieves them via dense vector similarity, but this approach overlooks the relationships between entities across chunks, making it insufficient for multi-hop reasoning tasks that require combining information from multiple sources.

Graph-based RAG methods~\cite{edge2024graphrag,guo2025lightrag,chen2025pathrag,gutierrez2025hipporag2} address this by structuring knowledge as graphs. However, they face three fundamental limitations:

\textbf{(1) Binary relation constraint.} Existing graph-based RAG methods represent knowledge as ordinary graphs where each edge connects exactly two entities. As shown by Luo et al.~\cite{luo2025hypergraphrag}, this binary representation is lossy---a fact involving $n \geq 3$ entities must be decomposed into multiple binary triples, causing information fragmentation and retrieval sparsity.

\textbf{(2) Expensive knowledge extraction.} HyperGraphRAG~\cite{luo2025hypergraphrag} addresses the first limitation by using hyperedges to represent $n$-ary relations, but relies on LLM-based extraction (\$0.006 per 1k tokens) for knowledge construction. Methods like LightRAG and PathRAG similarly require LLM calls for relation extraction during indexing.

\textbf{(3) Flat retrieval on structured knowledge.} Even methods that build rich graph structures often perform flat retrieval---top-$k$ similarity search without exploiting the graph topology for multi-step reasoning.

We propose \textbf{HyP-DLM} (Hypergraph Propagation with Dynamic Logic Modulation), a training-free framework that addresses all three limitations through the following contributions:

\begin{itemize}[leftmargin=*]
    \item \textbf{Cost-efficient hypergraph construction} via coreference resolution, semantic chunking, and local NER (GLiNER2), achieving near-zero indexing cost while maintaining information completeness.
    \item \textbf{DAG-guided matrix propagation} with degree-normalized incidence matrices, dynamic semantic modulation, and confidence-weighted synonym bridging based on a principled Bayesian formulation for uncertain entity links.
    \item \textbf{Multi-granular passage scoring} that combines hyperedge-level semantic relevance with entity-level propagation activation, replacing complex PPR-based ranking with a simple yet effective composite scoring mechanism.
\end{itemize}


%% ============================================================
%% 2. RELATED WORK
%% ============================================================
\section{Related Work}

\textbf{Graph-based RAG.} GraphRAG~\cite{edge2024graphrag} constructs knowledge graphs and retrieves community summaries. LightRAG~\cite{guo2025lightrag} improves efficiency via graph indexing. PathRAG~\cite{chen2025pathrag} refines retrieval with path pruning. HippoRAG2~\cite{gutierrez2025hipporag2} uses Personalized PageRank on binary graphs. All are constrained to binary relations.

\textbf{Hypergraph-based RAG.} HyperGraphRAG~\cite{luo2025hypergraphrag} is the first to use hyperedges for $n$-ary relation representation, achieving SOTA across multiple domains. However, its retrieval is flat (top-$k$ similarity) and its construction relies on LLM extraction.

\textbf{Multi-hop retrieval mechanisms.} BrowseNet~\cite{browsenet2026} decomposes queries into DAGs and traverses a graph-of-chunks in topological order. LinearRAG~\cite{linearrag2026} uses implicit semantic bridging via bipartite matrix propagation. MaGiX~\cite{hieu2025magix} proposes multi-granular scoring with composite chunk-attribute-triple scores for cross-lingual RAG. Our work synthesizes insights from all of these: DAG-guided reasoning structure, matrix-based propagation, multi-granular scoring, and confidence-weighted synonym bridging, unified within a hypergraph framework.


%% ============================================================
%% 3. METHOD
%% ============================================================
\section{Method: HyP-DLM}

HyP-DLM consists of three phases: (1) offline hypergraph construction via semantic chunking, (2) online DAG-guided propagation retrieval, and (3) trace-augmented generation.

%% ------------------------------------------------------------
%% 3.1 PHASE 1: INDEXING
%% ------------------------------------------------------------
\subsection{Phase 1: Cost-Efficient Hypergraph Construction}

Our key observation is that semantic chunks are natural hyperedges. A semantically coherent text fragment inherently captures an $n$-ary relational fact by describing a complete event or statement involving multiple entities. This eliminates the need for LLM-based relation extraction. The indexing pipeline consists of the following stages.

%% --- 3.1.1 Coreference Resolution ---
\subsubsection{Stage 1: Coreference Resolution}

Before chunking and NER, we perform document-level coreference resolution to replace pronouns and short mentions with their canonical forms. This is critical because semantic chunking produces smaller, focused chunks where local context may be insufficient to resolve references.

Given a document $d$, we apply a coreference resolution model (FastCoref with LingMessCoref) to produce a resolved document $d'$ where all mentions are replaced by their canonical form:
\begin{equation}
    d' = \text{CorefResolve}(d)
\end{equation}

The canonical selection follows a priority hierarchy: proper nouns $>$ common nouns $>$ pronouns, with longer spans preferred as tie-breakers. Post-processing extracts head nouns (e.g., ``rising star Lionel Messi'' $\to$ ``Lionel Messi'') and strips possessives and trailing punctuation.

\textbf{Example:}
\begin{quote}
    \textit{Input:} ``Messi scored. He celebrated with his teammates.'' \\
    \textit{Output:} ``Messi scored. Messi celebrated with Messi's teammates.''
\end{quote}

This ensures that downstream NER extracts ``Messi'' consistently rather than detecting ``He'' as a separate entity or missing the reference entirely.

%% --- 3.1.2 Semantic Chunking ---
\subsubsection{Stage 2: Semantic Chunking as Hyperedge Formation}

Given a resolved document $d'$, we split it into sentences $\{s_1, s_2, \ldots, s_L\}$ and compute sentence embeddings using a local model $f$ (all-mpnet-base-v2, 768 dimensions). We compute consecutive sentence similarity:
\begin{equation}
    \text{sim}(i) = \cos(f(s_i), f(s_{i+1})), \quad i \in [1, L-1]
\end{equation}

Breakpoints are inserted where $\text{sim}(i)$ falls below a percentile threshold of the similarity distribution within the document. Sentences between breakpoints form semantic chunks $C = \{c_1, c_2, \ldots, c_M\}$. Each chunk $c_j$ becomes a hyperedge $e_j$ in the knowledge hypergraph, with its text serving as the natural language description.

Chunk embeddings $f(c_j)$ are computed during this stage and \textbf{reused} throughout the pipeline (masking, modulation, context validation), avoiding redundant computation.

\begin{proposition}[Semantic Chunks as $N$-ary Facts]
A semantic chunk that describes a coherent event naturally connects $n \geq 2$ entities, preserving the $n$-ary relational structure without explicit extraction. Unlike atomic proposition decomposition (which reduces $n$-ary facts to unary statements), semantic chunks maintain the co-occurrence context of all involved entities.
\end{proposition}

%% --- 3.1.3 NER ---
\subsubsection{Stage 3: Entity Extraction via GLiNER2}

For each chunk $c_j$, we extract named entities using GLiNER2~\cite{zaratiana2025gliner2}, a unified schema-based information extraction model (205M parameters) that supports custom entity types through a schema interface:
\begin{equation}
    V_{c_j} = \text{GLiNER2}(c_j, \mathcal{L}) = \{(v_1^{\text{name}}, v_1^{\text{type}}), \ldots, (v_n^{\text{name}}, v_n^{\text{type}})\}
\end{equation}
where $\mathcal{L}$ is the set of target entity types (Person, Organization, Location, Event, Product, etc.), specified via GLiNER2's schema interface with optional descriptions for precision:
\begin{quote}
    \texttt{"person": "Full name of a real person, not pronouns or titles"}
\end{quote}

\textbf{Entity filtering.} We apply post-extraction filters to remove noise: entities with length $< 3$ characters (except known abbreviations), pure numeric strings, score patterns (e.g., ``3--1''), and entities whose NER type is not in the valid set.

After global deduplication across all chunks, we obtain the entity set $V = \{v_1, \ldots, v_N\}$.

\textbf{Why GLiNER2 over spaCy?} GLiNER2 allows custom entity types per domain without retraining, provides confidence scores for filtering, and supports relation extraction that can augment the synonym matrix. As a peer-reviewed model (EMNLP 2025), it offers both flexibility and credibility.

%% --- 3.1.4 Incidence Matrix ---
\subsubsection{Stage 4: Degree-Normalized Incidence Matrix}

We construct the raw incidence matrix $H \in \{0,1\}^{N \times M}$:
\begin{equation}
    H_{ij} = \begin{cases} 1 & \text{if entity } v_i \text{ appears in chunk } c_j \\ 0 & \text{otherwise} \end{cases}
\end{equation}

\textbf{Degree normalization.} Raw $H$ causes popular entities (e.g., ``United States'' appearing in 50 chunks) to dominate propagation through large self-loop scores in $HH^\top$. Following spectral hypergraph theory~\cite{zhou2006hypergraph}, we apply degree normalization:
\begin{equation}
    H_{\text{norm}} = D_v^{-1/2} \cdot H \cdot D_e^{-1/2}
\end{equation}
where:
\begin{itemize}[leftmargin=*]
    \item $D_v \in \mathbb{R}^{N \times N}$ is the diagonal entity degree matrix: $D_v[i,i] = \sum_{j=1}^{M} H_{ij}$ (number of hyperedges containing entity $v_i$).
    \item $D_e \in \mathbb{R}^{M \times M}$ is the diagonal hyperedge size matrix: $D_e[j,j] = \sum_{i=1}^{N} H_{ij}$ (number of entities in hyperedge $e_j$).
\end{itemize}

This normalization has two effects: (1) $D_v^{-1/2}$ penalizes high-degree entities, preventing popular entities from dominating the state vector; (2) $D_e^{-1/2}$ penalizes large hyperedges, giving more weight to co-occurrences in focused chunks (few entities) over broad chunks (many entities).

\textbf{Self-loop analysis.} After normalization, the self-loop score for entity $v_i$ is:
\begin{equation}
    (H_{\text{norm}} H_{\text{norm}}^\top)[i,i] = \sum_{j=1}^{M} \frac{H_{ij}}{D_v[i,i] \cdot D_e[j,j]}
\end{equation}
For entity ``United States'' (degree 50) and entity ``Rosario'' (degree 2), both self-loop scores are approximately 1.0 rather than 50 vs.\ 2, achieving balance.

%% --- 3.1.5 Synonym Matrix ---
\subsubsection{Stage 5: Confidence-Weighted Synonym Matrix}
\label{sec:synonym}

To handle entity name variations across chunks (e.g., ``Messi'' vs.\ ``Lionel Messi'', ``Bill Gates'' vs.\ ``William Henry Gates III''), we construct a synonym matrix $S_{\text{conf}} \in \mathbb{R}^{N \times N}$ with continuous confidence scores rather than binary links.

\textbf{Motivation: Bayesian propagation on uncertain graphs.} Entity-hyperedge links (captured by $H$) are deterministic---``Messi'' either appears in a chunk or not. In contrast, synonym links are inherently uncertain---``Messi'' and ``Lionel Messi'' are probably the same entity, but ``John Smith'' (scientist) and ``John Smith'' (athlete) may not be. Treating uncertain links as binary (present/absent) loses information about link reliability. We instead model synonym links probabilistically.

\begin{definition}[Confidence-Weighted Synonym Matrix]
Let $V = \{v_1, \ldots, v_N\}$ be the entity set. For each pair $(v_i, v_j)$, define:
\begin{equation}
    S_{\text{conf}}[i,j] = \begin{cases}
        \phi(\text{rrf}(v_i, v_j),\ \text{ctx}(v_i, v_j)) & \text{if pair is a candidate} \\
        0 & \text{otherwise}
    \end{cases}
\end{equation}
where $\phi$ is a confidence function, $\text{rrf}$ is the hybrid retrieval score, and $\text{ctx}$ is the context similarity score. The diagonal $S_{\text{conf}}[i,i] = 1$.
\end{definition}

\textbf{Candidate generation pipeline.} We use an 8-stage pipeline to identify synonym candidates efficiently:

\begin{enumerate}[leftmargin=*]
    \item \textbf{BM25 sparse encoding:} Each entity name is encoded using TF-IDF with character $n$-grams (bigram to 4-gram), capturing lexical overlap (e.g., ``FC Barcelona'' $\leftrightarrow$ ``Barcelona'' via shared character sequences).
    
    \item \textbf{Dense encoding:} Each entity name is encoded using all-mpnet-base-v2 (768 dimensions), capturing semantic similarity (e.g., ``Bar\c{c}a'' $\leftrightarrow$ ``FC Barcelona'' despite lexical difference).
    
    \item \textbf{Hybrid search via Qdrant:} For each entity, we query across documents using Reciprocal Rank Fusion (RRF) of sparse and dense results:
    \begin{equation}
        \text{rrf}(v_i, v_j) = \sum_{\ell \in \{\text{sparse}, \text{dense}\}} \frac{1}{k + \text{rank}_\ell(v_j)}
    \end{equation}
    where $k = 60$. Cross-document candidates only (no self-match).
    
    \item \textbf{Context validation (AND gate):} For each candidate pair $(v_i, v_j)$, compute context similarity using chunk embeddings:
    \begin{equation}
        \text{ctx}(v_i, v_j) = \max_{c_a \in \text{chunks}(v_i),\ c_b \in \text{chunks}(v_j)} \cos(f(c_a), f(c_b))
    \end{equation}
    Using MAX (not AVG): a single shared context is sufficient to confirm a link. Both $\text{rrf}$ and $\text{ctx}$ must independently exceed their respective thresholds for a pair to be retained.
\end{enumerate}

\textbf{Confidence function.} We define:
\begin{equation}
    \phi(\text{rrf}, \text{ctx}) = \text{rrf} \times \text{ctx}
\end{equation}
This product naturally attenuates uncertain links: high-confidence pairs (both scores high) receive strong weights, while borderline pairs receive weak weights.

\begin{theorem}[Probabilistic Interpretation]
\label{thm:bayesian}
If $S_{\text{conf}}[i,j]$ is interpreted as $P(v_i \equiv v_j \mid \text{evidence})$, i.e., the probability that entities $v_i$ and $v_j$ refer to the same real-world entity, then the propagation update:
\begin{equation}
    s_{\text{new}}[j] \mathrel{+}= S_{\text{conf}}[i,j] \times s[i]
\end{equation}
is the expected activation of $v_j$ under identity uncertainty:
\begin{equation}
    \mathbb{E}[s[j] \mid \text{uncertainty}] = P(v_i \equiv v_j) \cdot s[i] + P(v_i \not\equiv v_j) \cdot 0 = S_{\text{conf}}[i,j] \cdot s[i]
\end{equation}
\end{theorem}

\begin{corollary}
Binary synonym matrices ($S[i,j] \in \{0, 1\}$) are a special case when $P \in \{0, 1\}$, losing all information about link confidence. The continuous $S_{\text{conf}}$ enables graceful degradation when synonym links are uncertain.
\end{corollary}

%% --- 3.1.6 Semantic Masking ---
\subsubsection{Stage 6: Multi-Label Semantic Masking}
\label{sec:masking}

To reduce the search space during propagation, we cluster hyperedge embeddings and use cluster centroids as a coarse topic filter. Unlike standard single-label clustering where each chunk belongs to exactly one cluster, we employ \textbf{multi-label assignment} to avoid cutting cross-topic propagation paths.

\textbf{Problem with single-label clustering.} Consider a chunk ``Messi scored the winning goal for Barcelona in the Champions League.'' Under single-label HDBSCAN, this chunk is assigned to exactly one cluster (e.g., ``football''). A query about Messi's biography would activate the ``biography'' cluster mask, completely blocking this chunk from propagation---even though it contains the critical entity ``Messi'' needed for cross-topic bridging.

\textbf{Multi-centroid nearest neighbor (proposed).} We retain HDBSCAN for cluster discovery but add a multi-label post-processing step:

\begin{enumerate}[leftmargin=*]
    \item Run HDBSCAN on chunk embeddings $\{f(c_j)\}_{j=1}^M$ to obtain $K_{\text{auto}}$ clusters with centroids $\{\mu_k\}$ and noise labels.
    \item For each chunk $c_j$, compute similarity to all centroids:
    \begin{equation}
        \text{sim}_k(c_j) = \cos(f(c_j), \mu_k) \quad \forall k \in [1, K_{\text{auto}}]
    \end{equation}
    \item Assign chunk $c_j$ to multiple clusters via threshold:
    \begin{equation}
        \text{labels}(c_j) = \{k : \text{sim}_k(c_j) > \tau_{\text{multi}}\}
    \end{equation}
\end{enumerate}

During retrieval, given a guidance vector $g_i$ for sub-question $q_i$, the soft mask is:
\begin{equation}
    \text{mask}_j = \mathbbm{1}\left[\exists\, k \in \text{labels}(c_j) : k \in \text{Top-}P(\cos(g_i, \mu_k))\right] \lor \mathbbm{1}[\ell_j = -1]
\end{equation}
where Top-$P$ selects the $P$ clusters most similar to $g_i$, and noise points ($\ell_j = -1$) are always included as they may contain rare bridging facts.

\textbf{Returning to the Messi example:} The chunk is now assigned to both ``football'' and ``biography'' clusters. A biography-focused query activates both clusters, and the chunk survives masking.


%% ------------------------------------------------------------
%% 3.2 PHASE 2: RETRIEVAL
%% ------------------------------------------------------------
\subsection{Phase 2: DAG-Guided Propagation Retrieval}

%% --- 3.2.1 Query Decomposition ---
\subsubsection{Query Decomposition into DAG}

Unlike prior work that decomposes queries into linear chains of sub-questions, we decompose into a Directed Acyclic Graph (DAG) $\mathcal{Q} = (V_\mathcal{Q}, E_\mathcal{Q})$ to handle diverse reasoning structures:

\begin{itemize}[leftmargin=*]
    \item \textbf{Chain} (bridge-type): $q_1 \to q_2 \to q_3$
    \item \textbf{Parallel} (comparison-type): $q_1 \| q_2 \to q_3$
    \item \textbf{Tree} (mixed): $q_1 \to q_2,\ q_1 \to q_3,\ q_2 \wedge q_3 \to q_4$
\end{itemize}

We use a single LLM call with structured JSON output to decompose $q$ into sub-questions with dependency annotations. Each sub-question $q_i \in V_\mathcal{Q}$ has:
\begin{itemize}[leftmargin=*]
    \item A natural language question with \texttt{<ANS-$q_j$>} placeholders for dependencies.
    \item A guidance vector $g_i = f(q_i^{\text{resolved}})$ computed after placeholder replacement.
    \item A dependency set $\text{deps}(q_i) \subseteq V_\mathcal{Q}$.
\end{itemize}

Sub-questions are processed in topological order, ensuring all dependencies are resolved before a sub-question is processed. For simple queries, the LLM naturally returns a single-node DAG, making a separate routing mechanism unnecessary.

\textbf{Placeholder resolution} is a simple string operation (no additional LLM call):
\begin{equation}
    q_i^{\text{resolved}} = q_i^{\text{template}}.\text{replace}(\texttt{<ANS-}q_j\texttt{>},\ \text{answer}(q_j))
\end{equation}

%% --- 3.2.2 State Initialization ---
\subsubsection{State Initialization}

For each sub-question $q_i$ (after resolution), we initialize the state vector $\mathbf{s}_0 \in \mathbb{R}^N$ as follows.

\textbf{Step 1: NER-based seeding.} We extract entities from the resolved sub-question and match them to graph entities via embedding similarity:
\begin{equation}
    \mathbf{s}_0[k] = \max_{e \in \text{NER}(q_i^{\text{resolved}})} \cos(f(e), f(v_k)) \cdot \mathbbm{1}[\cos > \tau_{\text{init}}]
\end{equation}

\textbf{Fallback:} If NER extracts zero entities (e.g., the sub-question contains no named entities), we fall back to direct query-entity matching:
\begin{equation}
    \mathbf{s}_0[k] = \max(0,\ \cos(f(q_i^{\text{resolved}}), f(v_k))) \cdot \mathbbm{1}[\cos > \tau_{\text{init}}]
\end{equation}

\textbf{Step 2: Parent seeding.} For non-root sub-questions (those with dependencies), we additionally seed from parent activations:
\begin{equation}
    \mathbf{s}_0 = \max(\mathbf{s}_0^{\text{NER}},\ \alpha_{\text{parent}} \cdot \mathbf{s}_{\text{parent}}^{\text{final}})
\end{equation}
where $\mathbf{s}_{\text{parent}}^{\text{final}}$ is the converged state vector of the parent sub-question, $\alpha_{\text{parent}} = 0.5$ scales down parent signal to prevent domination, and $\max$ is element-wise.

%% --- 3.2.3 Propagation ---
\subsubsection{DAG-Guided Matrix Propagation}

This is the core retrieval algorithm. For each sub-question $q_i$ in topological order, we execute a propagation loop on the hypergraph guided by the sub-question's guidance vector.

\textbf{Phase A --- Semantic Masking (coarse filter).} Apply the multi-label masking strategy (Section~\ref{sec:masking}):
\begin{equation}
    \text{mask}_j = \mathbbm{1}\left[\exists\, k \in \text{labels}(c_j) : k \in \text{Top-}P(\cos(g_i, \mu_k))\right] \lor \mathbbm{1}[\ell_j = -1]
\end{equation}

\textbf{Phase B --- Dynamic Modulation (fine filter).} Compute the diagonal modulation matrix $D_i \in \mathbb{R}^{M \times M}$, which weights each hyperedge by its relevance to the current sub-question:
\begin{equation}
    D_i = \text{diag}\Big(\max\big(0,\ \cos(g_i, f(c_j))\big) \cdot \text{mask}_j\Big)_{j=1}^{M}
\end{equation}

The $\max(0, \cdot)$ (ReLU) ensures non-negative propagation. Cosine similarity is chosen over dot product (scale-invariant across embedding dimensions) and over softmax (which would force probability mass redistribution among surviving hyperedges, inflating irrelevant ones). Each $D_i[j,j] \in [0, 1]$:
\begin{itemize}[leftmargin=*]
    \item $= 0$ if hyperedge $j$ is masked out or semantically opposed to the query.
    \item $\approx 0$ if hyperedge $j$ is weakly related.
    \item $\approx 1$ if hyperedge $j$ is highly relevant to the sub-question.
\end{itemize}

\textbf{Important:} $D_i$ is fixed per sub-question (subscript $i$, not step $t$). The ``dynamic'' in Dynamic Logic Modulation refers to $D$ changing across sub-questions in the DAG ($D_1, D_2, \ldots$), not across propagation steps. This is critical for convergence (Section~\ref{sec:convergence}).

\textbf{Phase C --- Propagation Matrix (precomputed once per sub-question).}
\begin{equation}
\label{eq:propagation_matrix}
    A_i = H_{\text{norm}} \cdot D_i \cdot H_{\text{norm}}^\top + S_{\text{conf}}
\end{equation}
where:
\begin{itemize}[leftmargin=*]
    \item $H_{\text{norm}} \cdot D_i \cdot H_{\text{norm}}^\top \in \mathbb{R}^{N \times N}$: degree-normalized entity-to-entity propagation through modulated hyperedges. Signal flows entity $\to$ hyperedge $\to$ entity, with each hyperedge weighted by its relevance to $q_i$.
    \item $S_{\text{conf}} \in \mathbb{R}^{N \times N}$: confidence-weighted synonym bridging. Signal flows entity $\to$ synonym entity, with strength proportional to link confidence.
\end{itemize}

The two terms encode complementary propagation paths:
\begin{align}
    &\text{Path 1 (topology):} \quad v_a \xrightarrow{H_{\text{norm}}^\top} e_j \xrightarrow{D_i} e_j \xrightarrow{H_{\text{norm}}} v_b \\
    &\text{Path 2 (synonym):} \quad v_a \xrightarrow{S_{\text{conf}}} v_b
\end{align}

Note the absence of a separate weight $w_{\text{syn}}$ for the synonym term: $S_{\text{conf}}$ already encodes confidence levels (Section~\ref{sec:synonym}), making an additional scalar weight unnecessary.

\textbf{Phase D --- Damped PPR State Update.}
\begin{equation}
\label{eq:ppr}
    \mathbf{s}_t = (1 - \alpha) \cdot \mathbf{s}_0 + \alpha \cdot \frac{A_i \mathbf{s}_{t-1}}{\|A_i \mathbf{s}_{t-1}\|_\infty}
\end{equation}
where $\alpha \in (0, 1)$ is the damping factor (default 0.85). The teleport term $(1 - \alpha) \cdot \mathbf{s}_0$ continuously re-injects the initial query signal, preventing the propagated state from drifting away from the query intent. This is equivalent to Personalized PageRank on the hypergraph with personalization vector $\mathbf{s}_0$.

\textbf{Phase E --- Dynamic Pruning.}
\begin{equation}
    \mathbf{s}_t[k] = \begin{cases} \mathbf{s}_t[k] & \text{if } \mathbf{s}_t[k] > \delta \\ 0 & \text{otherwise} \end{cases}
\end{equation}
where $\delta = 0.05$. This keeps the state vector sparse, removing noise from distant, weakly activated entities and maintaining computational efficiency.

\textbf{Phase F --- Convergence Check.} The loop terminates when:
\begin{equation}
    \frac{\|\mathbf{s}_t - \mathbf{s}_{t-1}\|_2}{\|\mathbf{s}_{t-1}\|_2 + \epsilon} < \epsilon_{\text{conv}}
\end{equation}
or when $t = T$ (default $T = 6$, $\epsilon_{\text{conv}} = 0.01$, $\epsilon = 10^{-10}$).

%% --- 3.2.4 Multi-Granular Scoring ---
\subsubsection{Multi-Granular Passage Scoring}
\label{sec:scoring}

After propagation for each sub-question $q_i$, we score and rank chunks using a composite mechanism inspired by MaGiX~\cite{hieu2025magix}, combining two granularity levels.

\textbf{Hyperedge score (chunk-level relevance):} How relevant is the chunk to the sub-question?
\begin{equation}
    s_{\text{hyper}}(c_j, q_i) = D_i[j,j] = \max(0,\ \cos(g_i, f(c_j))) \cdot \text{mask}_j
\end{equation}

\textbf{Entity score (propagation-level activation):} How strongly are the entities in this chunk activated by the propagation?
\begin{equation}
    s_{\text{entity}}(c_j, q_i) = \max_{v_k \in \text{entities}(c_j)} \mathbf{s}_{\text{final}}[k]
\end{equation}
We use $\max$ (not mean) because a chunk is relevant if it contains \textit{at least one} highly activated entity.

\textbf{Composite score per sub-question:}
\begin{equation}
    s_{\text{comp}}(c_j, q_i) = w_h \cdot \hat{s}_{\text{hyper}}(c_j, q_i) + w_e \cdot \hat{s}_{\text{entity}}(c_j, q_i)
\end{equation}
where $\hat{s}$ denotes min-max normalized scores within each sub-question, and $w_h$, $w_e$ are weighting coefficients.

\textbf{Aggregation across sub-questions.} For each sub-question $q_i$, we retain the top-$K$ chunks by composite score. After processing all sub-questions in the DAG, we merge:
\begin{equation}
    s_{\text{final}}(c_j) = \max_{q_i \in V_\mathcal{Q}} s_{\text{comp}}(c_j, q_i)
\end{equation}
If a chunk appears in the top-$K$ of multiple sub-questions, its final score is the maximum across all sub-questions. The top-$N$ chunks by $s_{\text{final}}$ form the retrieved set $\mathcal{K}_{\text{graph}}$.

\textbf{Why two granularities?} Using only hyperedge scores captures semantic relevance but ignores whether the chunk lies on the reasoning path. Using only entity scores captures propagation activation but may be noisy from popular entities. Combining both ensures a chunk is \textit{both} semantically relevant \textit{and} on the reasoning path.

%% --- 3.2.5 Intermediate Answers ---
\subsubsection{Intermediate Answer Generation}

After scoring for each non-leaf sub-question $q_i$, we extract the top-$K_e$ entities and their associated top chunks. We generate a brief intermediate answer using the LLM, which is substituted into dependent sub-questions via \texttt{<ANS-$q_i$>} replacement. This intermediate answer also becomes part of the reasoning trace for final generation.


%% ------------------------------------------------------------
%% 3.3 PHASE 3: GENERATION
%% ------------------------------------------------------------
\subsection{Phase 3: Trace-Augmented Generation}

Unlike standard RAG which feeds only retrieved chunks to the LLM, we augment the generation prompt with the full reasoning trace from the DAG:

\begin{enumerate}[leftmargin=*]
    \item \textbf{Reasoning path:} The DAG structure with each sub-question and its intermediate answer, showing the multi-hop reasoning chain explicitly.
    \item \textbf{Key entities per hop:} The top activated entities for each sub-question, highlighting which entities were discovered at each reasoning step.
    \item \textbf{Ranked evidence chunks:} The top-$N$ chunks from the multi-granular scoring, ordered by composite score.
\end{enumerate}

The prompt structure:
\begin{quote}
\texttt{--- Reasoning Path ---} \\
\texttt{Sub-question 1: [question] $\to$ Answer: [intermediate answer]} \\
\texttt{Key entities: [entity list]} \\
\texttt{Sub-question 2: [resolved question] $\to$ Answer: [to be determined]} \\
\texttt{Key entities: [entity list]} \\[6pt]
\texttt{--- Evidence (ranked by relevance) ---} \\
\texttt{[1] (score: 0.92) "chunk text..."} \\
\texttt{[2] (score: 0.85) "chunk text..."} \\[6pt]
\texttt{--- Question ---} \\
\texttt{[original query]} \\[6pt]
\texttt{--- Instruction ---} \\
\texttt{Based on the reasoning path and evidence above, answer the question.} \\
\texttt{Reason inside <think>...</think>, answer inside <answer>...</answer>.}
\end{quote}

\textbf{Motivation:} The DAG decomposition and intermediate answers have already been computed (costing LLM calls). Not including them in the final prompt wastes this computation. By providing the reasoning trace, the LLM can verify and synthesize rather than re-discover the multi-hop chain, improving both accuracy and consistency.


%% ============================================================
%% 4. THEORETICAL ANALYSIS
%% ============================================================
\section{Theoretical Analysis}

%% --- 4.1 ---
\subsection{Semantic Chunks Preserve $N$-ary Information}

\begin{theorem}[Semantic Chunk Completeness]
Let $\mathcal{F} = \{v_1, \ldots, v_n\}$ be an $n$-ary fact expressed in a document. If all entities in $\mathcal{F}$ co-occur within a single semantic chunk $c_j$ (which holds when the fact is expressed in a coherent passage), then the hyperedge $e_j = c_j$ preserves the complete $n$-ary relation with zero information loss.
\end{theorem}

\begin{proof}[Proof sketch]
Since $c_j$ contains the full natural language description of the fact, and $V_{c_j} \supseteq \mathcal{F}$ (NER extracts all named entities), the hyperedge $(e_j, V_{c_j})$ connects all entities in $\mathcal{F}$ through a single structure. By the information-theoretic argument of~\cite{luo2025hypergraphrag}, this hypergraph representation has conditional entropy $H(X | \phi_H(X)) = 0$, preserving all information.
\end{proof}

\textbf{When does this fail?} If a fact spans multiple sentences that are split into different chunks, information is fragmented. We mitigate this with (a) coreference resolution ensuring entity names are explicit, (b) synonym matrix $S_{\text{conf}}$ connecting entities across chunks, and (c) multi-step propagation that naturally bridges across chunks.

%% --- 4.2 ---
\subsection{DAG Propagation Covers All Reasoning Patterns}

\begin{theorem}
Any multi-hop question with $k$ reasoning steps can be represented as a DAG with at most $k$ nodes. Chain, comparison, bridge, and tree reasoning patterns are all special cases of DAGs.
\end{theorem}

This follows from the observation that reasoning dependencies in natural questions are acyclic (an answer cannot depend on itself). BrowseNet~\cite{browsenet2026} empirically validates that all gold decompositions in HotpotQA, 2WikiMQA, and MuSiQue are DAGs.

%% --- 4.3 ---
\subsection{Propagation Convergence}
\label{sec:convergence}

\begin{theorem}[Damped PPR Convergence on Hypergraph]
\label{thm:convergence}
Let $A_i = H_{\text{norm}} D_i H_{\text{norm}}^\top + S_{\text{conf}}$ be the fixed propagation matrix for sub-question $q_i$. The damped PPR update (Eq.~\ref{eq:ppr}) converges to a unique fixed point $\mathbf{s}_i^*$ at geometric rate $O(\alpha^t)$.
\end{theorem}

\begin{proof}
Define the operator $T(\mathbf{s}) = (1-\alpha) \cdot \mathbf{s}_0 + \alpha \cdot \bar{A}_i \mathbf{s}$, where $\bar{A}_i$ denotes $A_i$ followed by $\ell_\infty$-normalization. Since $\bar{A}_i$ maps vectors into $[0, 1]^N$ (after normalization), $T$ maps into a convex compact set. Moreover:
\begin{equation}
    \|T(\mathbf{s}) - T(\mathbf{s}')\| = \alpha \|\bar{A}_i(\mathbf{s} - \mathbf{s}')\| \leq \alpha \|\mathbf{s} - \mathbf{s}'\|
\end{equation}
Since $0 < \alpha < 1$, $T$ is a contraction mapping. By the Banach fixed-point theorem, there exists a unique fixed point $\mathbf{s}_i^* = T(\mathbf{s}_i^*)$, and:
\begin{equation}
    \|\mathbf{s}_t - \mathbf{s}_i^*\| \leq \alpha^t \|\mathbf{s}_0 - \mathbf{s}_i^*\|
\end{equation}
For $\alpha = 0.85$, convergence to $\epsilon$-accuracy requires $t \geq \lceil \log(\epsilon) / \log(\alpha) \rceil$ steps. In practice, we observe convergence within $T = 6$ steps.
\end{proof}

\textbf{Role of pruning.} Dynamic pruning ($\mathbf{s}_t[k] = 0$ if $\mathbf{s}_t[k] < \delta$) is a projection onto the set $\{\mathbf{s} : \mathbf{s}[k] \geq \delta \text{ or } \mathbf{s}[k] = 0\}$. Since the teleport term $(1-\alpha) \cdot \mathbf{s}_0$ re-injects the initial signal at each step, entities with non-zero $\mathbf{s}_0[k]$ are restored if their initial activation exceeds $\delta / (1 - \alpha)$.

%% --- 4.4 ---
\subsection{Confidence-Attenuated Propagation}

\begin{theorem}[Optimal Propagation under Identity Uncertainty]
Let $G^* = (V, E_H, E_S)$ be the extended hypergraph with certain edges $E_H$ (entity-hyperedge) and uncertain edges $E_S$ (synonym links). If $S_{\text{conf}}[i,j] = P(v_i \equiv v_j \mid \text{evidence})$, then the propagation matrix $A_i = H_{\text{norm}} D_i H_{\text{norm}}^\top + S_{\text{conf}}$ yields the maximum likelihood estimator for entity activation under identity uncertainty:
\begin{equation}
    \mathbb{E}[\mathbf{s}_t \mid \text{identity uncertainty}] = A_i \cdot \mathbf{s}_{t-1}
\end{equation}
\end{theorem}

\begin{proof}[Proof sketch]
For each entity pair $(v_i, v_j)$ connected by a synonym edge, the signal transmitted is:
\begin{equation}
    \mathbb{E}[s_{\text{new}}[j]] = P(v_i \equiv v_j) \cdot s[i] + P(v_i \not\equiv v_j) \cdot 0 = S_{\text{conf}}[i,j] \cdot s[i]
\end{equation}
For certain edges ($H$), the probability is 1, reducing to standard propagation. Summing over both edge types yields $A_i \cdot \mathbf{s}_{t-1}$.
\end{proof}

%% --- 4.5 ---
\subsection{Propagation Complexity}

Each propagation step involves one sparse matrix-vector multiplication with $A_i$. Let $\text{nnz}(H)$ denote the number of non-zeros in $H$:
\begin{equation}
    \text{Cost per step} = O(\text{nnz}(A_i)) = O(\text{nnz}(H)^2 / M + \text{nnz}(S_{\text{conf}}))
\end{equation}

In practice, with masking reducing $M$ to $\sim M' / K$ effective hyperedges and typical hypergraph density ($\sim$3--5 entities per hyperedge), each step takes $< 1$ms with SciPy sparse operations. The total cost per sub-question is $O(T \cdot \text{nnz}(A_i))$ where $T$ is the number of propagation steps (typically 3--5 with early stopping). For a DAG with $|V_\mathcal{Q}|$ sub-questions, total propagation cost is $O(|V_\mathcal{Q}| \cdot T \cdot \text{nnz}(A_i))$, dominated by LLM calls for query decomposition and intermediate answer generation.


%% ============================================================
%% 5. EXPERIMENTS
%% ============================================================
\section{Experiments}

\subsection{Experimental Setup}

\textbf{Datasets.} We evaluate on three multi-hop QA benchmarks: HotpotQA~\cite{yang2018hotpotqa} (2-hop, 7,405 dev), 2WikiMultiHopQA~\cite{ho2020constructing} (2--4 hop, 12,576 dev), and MuSiQue~\cite{trivedi2022musique} (2--4 hop, 2,417 dev).

\textbf{Baselines.} We compare against: NaiveGeneration, StandardRAG, GraphRAG~\cite{edge2024graphrag}, LightRAG~\cite{guo2025lightrag}, PathRAG~\cite{chen2025pathrag}, HippoRAG2~\cite{gutierrez2025hipporag2}, HyperGraphRAG~\cite{luo2025hypergraphrag}, BrowseNet~\cite{browsenet2026}, and LinearRAG~\cite{linearrag2026}.

\textbf{Metrics.} F1 (word-level overlap), Retrieval Similarity (R-S, semantic cosine between retrieved and gold knowledge), and Generation Evaluation (G-E, 7-dimension LLM judge + F1).

\textbf{Implementation.} Embedding model: all-mpnet-base-v2 (768d). NER: GLiNER2 (gliner2-base-v1, 205M). Coreference: FastCoref (LingMessCoref). LLM: GPT-4o-mini for query decomposition and generation.

%% --- 5.2 Main Results ---
\subsection{Main Results}

\begin{table}[ht]
\centering
\caption{Performance comparison on multi-hop QA benchmarks. Bold indicates best, underline indicates second best. [To be filled after experiments.]}
\label{tab:main}
\begin{tabular}{l ccc ccc ccc}
\toprule
\multirow{2}{*}{Method} & \multicolumn{3}{c}{HotpotQA} & \multicolumn{3}{c}{2WikiMQA} & \multicolumn{3}{c}{MuSiQue} \\
\cmidrule(lr){2-4} \cmidrule(lr){5-7} \cmidrule(lr){8-10}
 & F1 & R-S & G-E & F1 & R-S & G-E & F1 & R-S & G-E \\
\midrule
NaiveGeneration & -- & -- & -- & -- & -- & -- & -- & -- & -- \\
StandardRAG & -- & -- & -- & -- & -- & -- & -- & -- & -- \\
GraphRAG & -- & -- & -- & -- & -- & -- & -- & -- & -- \\
LightRAG & -- & -- & -- & -- & -- & -- & -- & -- & -- \\
PathRAG & -- & -- & -- & -- & -- & -- & -- & -- & -- \\
HippoRAG2 & -- & -- & -- & -- & -- & -- & -- & -- & -- \\
HyperGraphRAG & -- & -- & -- & -- & -- & -- & -- & -- & -- \\
BrowseNet & -- & -- & -- & -- & -- & -- & -- & -- & -- \\
LinearRAG & -- & -- & -- & -- & -- & -- & -- & -- & -- \\
\midrule
HyP-DLM (Ours) & -- & -- & -- & -- & -- & -- & -- & -- & -- \\
\bottomrule
\end{tabular}
\end{table}

%% --- 5.3 Ablation ---
\subsection{Ablation Study: Core Components}

We ablate each component on HotpotQA to measure individual contributions.

\begin{table}[ht]
\centering
\caption{Component ablation study on HotpotQA. Each row removes one component from the full system. [To be filled after experiments.]}
\label{tab:ablation}
\begin{tabular}{l ccc}
\toprule
Variant & F1 & R-S & G-E \\
\midrule
HyP-DLM (Full) & -- & -- & -- \\
\midrule
w/o DAG decomposition (single query, no decompose) & -- & -- & -- \\
w/o Synonym Matrix $S_{\text{conf}}$ & -- & -- & -- \\
w/o Semantic Masking (mask$_j$ = 1 for all) & -- & -- & -- \\
w/o Dynamic Modulation $D_i$ ($D_i = I$) & -- & -- & -- \\
w/o Degree Normalization ($H$ raw) & -- & -- & -- \\
w/o Dynamic Pruning ($\delta = 0$) & -- & -- & -- \\
w/o DAG trace in generation (chunks only) & -- & -- & -- \\
\bottomrule
\end{tabular}
\end{table}

%% --- 5.4 Sensitivity ---
\subsection{Sensitivity Analysis}

We study the sensitivity of key hyperparameters on HotpotQA.

\begin{table}[ht]
\centering
\caption{Sensitivity of multi-label masking threshold $\tau_{\text{multi}}$ on HotpotQA. [To be filled.]}
\label{tab:sens_mask}
\begin{tabular}{l ccc cc}
\toprule
$\tau_{\text{multi}}$ & F1 & R-S & G-E & Mask \% & Lat. (ms) \\
\midrule
0.3 & -- & -- & -- & -- & -- \\
0.4 & -- & -- & -- & -- & -- \\
0.5 & -- & -- & -- & -- & -- \\
0.6 & -- & -- & -- & -- & -- \\
0.7 & -- & -- & -- & -- & -- \\
\bottomrule
\end{tabular}
\end{table}

\begin{table}[ht]
\centering
\caption{Sensitivity of damping factor $\alpha$ and pruning threshold $\delta$ on HotpotQA. [To be filled.]}
\label{tab:sens_prop}
\begin{tabular}{ll ccc cc}
\toprule
$\alpha$ & $\delta$ & F1 & R-S & G-E & Avg Hops & $|$Active$|$ \\
\midrule
0.50 & 0.05 & -- & -- & -- & -- & -- \\
0.70 & 0.05 & -- & -- & -- & -- & -- \\
0.85 & 0.05 & -- & -- & -- & -- & -- \\
0.90 & 0.05 & -- & -- & -- & -- & -- \\
0.95 & 0.05 & -- & -- & -- & -- & -- \\
\midrule
0.85 & 0.01 & -- & -- & -- & -- & -- \\
0.85 & 0.03 & -- & -- & -- & -- & -- \\
0.85 & 0.05 & -- & -- & -- & -- & -- \\
0.85 & 0.10 & -- & -- & -- & -- & -- \\
0.85 & 0.20 & -- & -- & -- & -- & -- \\
\bottomrule
\end{tabular}
\end{table}

\begin{table}[ht]
\centering
\caption{Sensitivity of composite scoring weights on HotpotQA. [To be filled.]}
\label{tab:sens_scoring}
\begin{tabular}{cc ccc}
\toprule
$w_h$ & $w_e$ & F1 & R-S & G-E \\
\midrule
1.0 & 0.3 & -- & -- & -- \\
1.0 & 0.5 & -- & -- & -- \\
0.5 & 0.5 & -- & -- & -- \\
0.3 & 1.0 & -- & -- & -- \\
\bottomrule
\end{tabular}
\end{table}

%% --- 5.5 Cost ---
\subsection{Cost Analysis}

\begin{table}[ht]
\centering
\caption{Indexing and generation cost comparison. [To be filled.]}
\label{tab:cost}
\begin{tabular}{l cc cc}
\toprule
\multirow{2}{*}{Method} & \multicolumn{2}{c}{Construction} & \multicolumn{2}{c}{Generation} \\
\cmidrule(lr){2-3} \cmidrule(lr){4-5}
 & Time/1kT & Cost/1kT & Time/Query & Cost/1kQ \\
\midrule
StandardRAG & 0 s & \$0 & -- s & -- \\
GraphRAG & 9.27 s & \$0.0058 & 0.22 s & \$1.84 \\
LightRAG & 5.17 s & \$0.0081 & 0.36 s & \$3.36 \\
HyperGraphRAG & 3.08 s & \$0.0063 & 0.26 s & \$3.18 \\
HyP-DLM (Ours) & -- s & $\sim$\$0 & -- s & -- \\
\bottomrule
\end{tabular}
\end{table}


%% ============================================================
%% 6. DISCUSSION
%% ============================================================
\section{Discussion}

\textbf{When does semantic chunking underperform LLM extraction?} When facts are expressed across long, complex paragraphs where entity co-occurrence within a single chunk is not guaranteed, LLM-based extraction may identify relations that span chunk boundaries. Our coreference resolution and multi-step propagation mitigate this but may not fully compensate for very long-range dependencies.

\textbf{Scalability.} The matrix propagation approach scales linearly with graph size (dominated by sparse matrix operations). For very large corpora ($> 1$M chunks), FAISS-based approximate nearest neighbor search can replace brute-force similarity computation in the masking phase.

\textbf{Limitations.} (1) Query decomposition requires one LLM call, adding latency. (2) Coreference resolution (FastCoref) is the computational bottleneck at indexing time ($\sim$3s/doc on CPU). (3) NER-based entity extraction may miss implicit entities not expressed as named entities.


%% ============================================================
%% 7. CONCLUSION
%% ============================================================
\section{Conclusion}

We presented HyP-DLM, a training-free hypergraph-based RAG framework that achieves cost-efficient knowledge construction via coreference resolution, semantic chunking, and local NER; effective multi-hop retrieval via DAG-guided matrix propagation with degree-normalized incidence matrices, confidence-weighted synonym bridging, and multi-label semantic masking; and robust passage ranking via multi-granular composite scoring. Our theoretical contributions include a Bayesian formulation for propagation on uncertain entity graphs (Theorem~\ref{thm:bayesian}) and convergence guarantees for damped PPR on hypergraphs (Theorem~\ref{thm:convergence}). Experiments demonstrate competitive or superior performance compared to state-of-the-art methods while maintaining near-zero indexing cost.


%% ============================================================
%% REFERENCES
%% ============================================================
\bibliographystyle{plain}
\begin{thebibliography}{20}

\bibitem{lewis2020rag}
P.~Lewis et al., ``Retrieval-augmented generation for knowledge-intensive NLP tasks,'' \textit{NeurIPS}, 2020.

\bibitem{edge2024graphrag}
D.~Edge et al., ``From local to global: A graph RAG approach to query-focused summarization,'' 2024.

\bibitem{guo2025lightrag}
Z.~Guo et al., ``LightRAG: Simple and fast retrieval-augmented generation,'' 2024.

\bibitem{chen2025pathrag}
B.~Chen et al., ``PathRAG: Pruning graph-based retrieval augmented generation with relational paths,'' 2025.

\bibitem{gutierrez2025hipporag2}
B.~J.~Guti\'errez et al., ``From RAG to memory: Non-parametric continual learning for large language models,'' 2025.

\bibitem{luo2025hypergraphrag}
H.~Luo et al., ``HyperGraphRAG: Retrieval-augmented generation via hypergraph-structured knowledge representation,'' \textit{NeurIPS}, 2025.

\bibitem{browsenet2026}
BrowseNet, ``Graph-based associative memory for contextual information retrieval,'' \textit{ICLR submission}, 2026.

\bibitem{linearrag2026}
LinearRAG, ``Linear retrieval-augmented generation via tri-graph semantic bridging,'' 2026.

\bibitem{hieu2025magix}
N.~M.~Hieu et al., ``MaGiX: A multi-granular adaptive graph intelligence framework for enhancing cross-lingual RAG,'' \textit{Findings of EMNLP}, 2025.

\bibitem{zaratiana2025gliner2}
U.~Zaratiana et al., ``GLiNER2: Schema-driven multi-task learning for structured information extraction,'' \textit{EMNLP System Demonstrations}, 2025.

\bibitem{zhou2006hypergraph}
D.~Zhou et al., ``Learning with hypergraphs: Clustering, classification, and embedding,'' \textit{NeurIPS}, 2006.

\bibitem{yang2018hotpotqa}
Z.~Yang et al., ``HotpotQA: A dataset for diverse, explainable multi-hop question answering,'' \textit{EMNLP}, 2018.

\bibitem{ho2020constructing}
X.~Ho et al., ``Constructing a multi-hop QA dataset for comprehensive evaluation of reasoning steps,'' \textit{COLING}, 2020.

\bibitem{trivedi2022musique}
H.~Trivedi et al., ``MuSiQue: Multihop questions via single hop question composition,'' \textit{TACL}, 2022.

\bibitem{campello2013hdbscan}
R.~J.~G.~B.~Campello, D.~Moulavi, and J.~Sander, ``Density-based clustering based on hierarchical density estimates,'' \textit{PAKDD}, 2013.

\bibitem{johnson2019faiss}
J.~Johnson, M.~Douze, and H.~J\'egou, ``Billion-scale similarity search with GPUs,'' \textit{IEEE Trans.\ Big Data}, 2019.

\end{thebibliography}


%% ============================================================
%% APPENDIX
%% ============================================================
\appendix

\section{Full Algorithm Pseudocode}

\begin{algorithm}[ht]
\caption{HyP-DLM: DAG-Guided Propagation Retrieval}
\label{alg:hypdlm}
\begin{algorithmic}[1]
\Require Query $q$, Hypergraph $(H_{\text{norm}}, S_{\text{conf}}, \{f(c_j)\}, \{\mu_k\})$
\Ensure Ranked passages $\mathcal{K}_{\text{graph}}$

\State \textbf{// Step 1: Query Decomposition}
\State $\mathcal{Q} = (V_\mathcal{Q}, E_\mathcal{Q}) \gets \text{LLM\_Decompose}(q)$ \Comment{Single LLM call, JSON output}

\State \textbf{// Step 2: Process sub-questions in topological order}
\State $\text{all\_chunks} \gets \{\}$
\For{$q_i \in \text{TopologicalSort}(V_\mathcal{Q})$}
    \State $q_i^{\text{res}} \gets \text{ResolvePlaceholders}(q_i)$ \Comment{String replace, no LLM}
    \State $g_i \gets f(q_i^{\text{res}})$ \Comment{Guidance vector}
    
    \State \textbf{// State initialization}
    \State $\text{ents} \gets \text{NER}(q_i^{\text{res}})$
    \If{$|\text{ents}| > 0$}
        \State $\mathbf{s}_0[k] \gets \max_{e \in \text{ents}} \cos(f(e), f(v_k)) \cdot \mathbbm{1}[\cos > \tau_{\text{init}}]$
    \Else
        \State $\mathbf{s}_0[k] \gets \max(0, \cos(g_i, f(v_k))) \cdot \mathbbm{1}[\cos > \tau_{\text{init}}]$ \Comment{Fallback}
    \EndIf
    \If{$\text{deps}(q_i) \neq \emptyset$}
        \State $\mathbf{s}_0 \gets \max(\mathbf{s}_0, \alpha_{\text{parent}} \cdot \mathbf{s}_{\text{parent}})$ \Comment{Parent seeding}
    \EndIf
    
    \State \textbf{// Compute mask and modulation (once per sub-question)}
    \State $\text{mask}_j \gets \mathbbm{1}[\exists\, k \in \text{labels}(c_j) : k \in \text{Top-}P(\cos(g_i, \mu_k))] \lor \mathbbm{1}[\ell_j = -1]$
    \State $D_i \gets \text{diag}(\max(0, \cos(g_i, f(c_j))) \cdot \text{mask}_j)_{j=1}^{M}$
    \State $A_i \gets H_{\text{norm}} \cdot D_i \cdot H_{\text{norm}}^\top + S_{\text{conf}}$
    
    \State \textbf{// Propagation loop}
    \For{$t = 1, \ldots, T$}
        \State $\mathbf{s}_t \gets (1 - \alpha) \cdot \mathbf{s}_0 + \alpha \cdot A_i \mathbf{s}_{t-1} / \|A_i \mathbf{s}_{t-1}\|_\infty$
        \State $\mathbf{s}_t[k] \gets 0$ if $\mathbf{s}_t[k] < \delta$ \Comment{Pruning}
        \If{$\|\mathbf{s}_t - \mathbf{s}_{t-1}\|_2 / (\|\mathbf{s}_{t-1}\|_2 + \epsilon) < \epsilon_{\text{conv}}$}
            \State \textbf{break} \Comment{Converged}
        \EndIf
    \EndFor
    
    \State \textbf{// Multi-granular scoring}
    \For{each chunk $c_j$}
        \State $s_{\text{hyper}}(c_j) \gets D_i[j,j]$
        \State $s_{\text{entity}}(c_j) \gets \max_{v_k \in \text{entities}(c_j)} \mathbf{s}_t[k]$
        \State $s_{\text{comp}}(c_j) \gets w_h \cdot \hat{s}_{\text{hyper}} + w_e \cdot \hat{s}_{\text{entity}}$ \Comment{Min-max normalized}
    \EndFor
    \State $\text{all\_chunks}[q_i] \gets \text{Top-}K(s_{\text{comp}})$
    
    \State \textbf{// Intermediate answer (non-leaf nodes)}
    \If{$q_i$ is not a leaf}
        \State $\text{answer}(q_i) \gets \text{LLM\_Answer}(q_i^{\text{res}}, \text{Top-}K \text{ chunks})$
    \EndIf
\EndFor

\State \textbf{// Step 3: Merge and rank}
\State $s_{\text{final}}(c_j) \gets \max_{q_i} s_{\text{comp}}(c_j, q_i)$
\State $\mathcal{K}_{\text{graph}} \gets \text{Top-}N(s_{\text{final}})$

\State \Return $\mathcal{K}_{\text{graph}}$
\end{algorithmic}
\end{algorithm}


\section{Prompt Templates}

\subsection{Query Decomposition Prompt}
\begin{verbatim}
You are a question decomposition expert. Given a 
complex question, break it down into simple single-hop 
sub-questions and their dependencies.

Output format (JSON):
{
  "nodes": [
    {"id": "q1", "question": "...", "depends_on": []},
    {"id": "q2", "question": "...", "depends_on": ["q1"]}
  ]
}

Rules:
- Each sub-question should be answerable from a single 
  passage.
- Use <ANS-qN> to reference the answer of sub-question qN.
- The DAG must be acyclic.
- For comparison questions, create parallel branches.
- For simple questions, return a single node.

Question: {query}
\end{verbatim}


\section{Hyperparameter Settings}

\begin{table}[ht]
\centering
\caption{Full hyperparameter settings for HyP-DLM.}
\label{tab:hyperparams}
\small
\begin{tabular}{llc}
\toprule
Phase & Parameter & Value \\
\midrule
\multirow{4}{*}{Coreference} 
    & Model & FastCoref (LingMessCoref) \\
    & Priority & Proper $>$ Common $>$ Pronoun \\
    & Tie-break & Longer span \\
    & Post-process & Head extraction + strip \\
\midrule
\multirow{4}{*}{Chunking}
    & Embedding model & all-mpnet-base-v2 (768d) \\
    & Breakpoint type & Percentile \\
    & Min chunk tokens & 50 \\
    & Max chunk tokens & 300 \\
\midrule
\multirow{3}{*}{NER}
    & Model & GLiNER2 (gliner2-base-v1) \\
    & Min entity length & 3 characters \\
    & Valid types & PER, ORG, GPE, LOC, ... \\
\midrule
\multirow{4}{*}{Synonym}
    & Sparse encoder & TF-IDF char $n$-gram (2--4) \\
    & Dense encoder & all-mpnet-base-v2 \\
    & RRF $k$ & 60 \\
    & Context validation & MAX cosine, AND gate \\
\midrule
\multirow{3}{*}{Masking}
    & Strategy & Multi-centroid NN \\
    & Base clustering & HDBSCAN (min\_cluster=15) \\
    & $\tau_{\text{multi}}$ & 0.5 \\
\midrule
\multirow{8}{*}{Propagation}
    & Damping $\alpha$ & 0.85 \\
    & Convergence $\epsilon_{\text{conv}}$ & 0.01 \\
    & Max hops $T$ & 6 \\
    & Top-$P$ clusters & 5 \\
    & Pruning threshold $\delta$ & 0.05 \\
    & Init threshold $\tau_{\text{init}}$ & 0.5 \\
    & Parent seeding $\alpha_{\text{parent}}$ & 0.5 \\
\midrule
\multirow{3}{*}{Scoring}
    & $w_h$ (hyperedge weight) & 1.0 \\
    & $w_e$ (entity weight) & 0.3 \\
    & Top-$K$ per sub-question & 30 \\
    & Top-$N$ final & 10 \\
\midrule
Generation & LLM model & GPT-4o-mini \\
\bottomrule
\end{tabular}
\end{table}


\end{document}